\chapter{Ranking, Recommendation, and Exposure}
\label{ch:ranking-recommendation}

\emph{Opening dilemma.} A course-support team launches a recommender that suggests five study resources whenever the early-warning model flags a student as CONCERNING. After six months they audit it. Worked-example videos --- the resource the team showed most often last semester --- now drive 80\% of recommendations. Practice problems and tutoring sign-ups, which a quieter cohort of students preferred, have nearly disappeared from the top-5 lists. The model is not broken in any conventional sense: its offline replay accuracy on last semester's logs is higher than ever. The replay just \emph{is} last semester's policy. The system has taught itself to recommend, more and more confidently, exactly what last semester's system already showed --- and it has erased the evidence that any other recommendation could have worked. That feedback loop is what this chapter is really about. Collaborative filtering, matrix factorization, and NDCG are the methods; exposure bias is the modeling dilemma they have to answer.

Earlier chapters built models that classify, regress, detect, and forecast. In each case, the model produces one prediction per input. But many real-world systems do not predict a single answer---they produce a \emph{ranked list}. A search engine ranks documents. A streaming service ranks movies. A course-support system ranks resources for a struggling student. A campus job board ranks opportunities for graduating seniors.

Ranking introduces problems classification and regression do not face. The items a system surfaces shape what users see, click, and ultimately choose. The model's outputs influence the data it will be trained on next---a feedback loop that can amplify biases, suppress diversity, and create blind spots. The design question shifts from ``is the prediction correct?'' to ``is the ranked list useful, fair, and safe?''

This chapter treats recommendation as a ranking problem. It introduces collaborative filtering and matrix factorization, addresses the challenges of implicit feedback and missing data, and examines how exposure bias corrupts evaluation and perpetuates inequality. The chapter artifact---the recommendation audit card---documents what the system exposes, suppresses, and amplifies.

A recommender system is not a neutral mirror of user preferences. It actively shapes what users experience. Every design decision---what to rank, which signals to use, how to handle missing data---decides what gets seen and what stays hidden.

\textbf{Exposure as allocation.} When evaluating a recommendation system, ask who or what is being exposed and who or what is being suppressed. Ranking allocates user attention, and that allocation has consequences for content creators, students, job seekers, and anyone whose opportunities depend on being visible.

\begin{practicalnote}
\textbf{Depth guide.} The core target is to understand ranking as exposure allocation and to recognize the feedback-loop risks that ordinary classifiers do not create. Full-scale recommender optimization, production counterfactual logging, and marketplace fairness are natural extensions for projects.
\end{practicalnote}

\section{Recommendation as Ranking}

A recommendation system predicts which items a user will find relevant and then presents them in a ranked order. This framing connects to several ideas already encountered in the book.

\textbf{Connection to classification.} At the item level, recommendation can sometimes be framed as predicting a relevance or utility score for a user-item pair. That score may come from a pointwise classifier or regressor, a pairwise preference model, or a listwise ranker. The ranked list is produced by sorting or otherwise selecting from candidate scores.

\textbf{Connection to retrieval.} Information retrieval---search engines, document retrieval---also produces ranked lists. The key difference is that search is query-driven (the user states what they want) while recommendation is profile-driven (the system infers what the user might want from past behavior). Evaluation metrics overlap: precision, recall, and normalized discounted cumulative gain (NDCG) appear in both settings.

\textbf{Connection to representation learning.} The most successful recommendation systems learn latent representations of users and items (Chapter~\ref{ch:representation-foundation-models}). A user is a vector that captures their preferences; an item is a vector that captures its attributes. Recommendation becomes a learned scoring problem in the latent space: recommend items whose vectors are most compatible with the user's vector under a score such as a dot product or cosine similarity.

\begin{example}[Course Resource Recommendation]
A university's course-support system has a library of 500 resources: lecture recordings, worked examples, practice problems, office-hour transcripts, and tutoring sessions. When the early-warning model flags a student as CONCERNING, the system should recommend the 5 most helpful resources. This is a ranking problem: score all 500 resources by predicted helpfulness for this student and present the top 5. The scoring function must account for the student's specific weaknesses (revealed by their submission patterns), the resource's content, and what similar students found helpful in the past.
\end{example}

\section{Collaborative Filtering}
\label{sec:collab-filter}

\subsection{The Core Idea}

Collaborative filtering is the foundational approach to recommendation. The idea is simple: recommend items similar users liked.

\textbf{User-based collaborative filtering.} To recommend items for user $u$:
\begin{enumerate}[leftmargin=*]
\item Find users who are similar to $u$ (based on their past ratings or interactions).
\item Recommend items that those similar users liked but $u$ has not yet seen.
\end{enumerate}

\textbf{Item-based collaborative filtering.} To recommend items for user $u$:
\begin{enumerate}[leftmargin=*]
\item For each item $u$ has liked, find similar items (based on which users liked them).
\item Recommend the most similar items that $u$ has not yet seen.
\end{enumerate}

Both approaches require defining ``similarity.'' Common choices are cosine similarity (treating each user's ratings as a vector) and Pearson correlation (which adjusts for different rating scales across users).

\subsection{The Cold-Start Problem}

Collaborative filtering fails when there is insufficient interaction data:

\textbf{New users.} A student who just enrolled has no interaction history. The system cannot find similar users because nothing is yet comparable.

\textbf{New items.} A newly uploaded resource has no ratings. The system cannot recommend it because no users have interacted with it.

\textbf{Sparse interactions.} Even for established users, most user-item pairs have no interaction. A university with 1,000 students and 500 resources has 500,000 possible interactions, yet each student may have accessed only 20 resources. The interaction matrix is 96\% empty.

The cold-start problem is not a technical nuisance---it is a design problem. New items that are never recommended will never collect interactions, so they will keep going unrecommended. This creates a rich-get-richer dynamic: popular items accumulate more data and become even more popular, while new or niche items stay permanently invisible. Addressing cold start requires deliberate exploration---sometimes recommending items the system is uncertain about, not just items it is confident the user will like.

\begin{example}[Cold Start on a New Resource]
The course-support recommender has just added a new resource---``Quiz: Linear Algebra Review''---with zero interaction history. Three strategies the team could use, and what each one quietly bets on:

\emph{1. Popularity fallback.} Rank items by total clicks across all students and recommend the top ones. The new quiz is never shown because it has zero clicks, so it never accumulates evidence, so it is never shown. The state is self-reinforcing and the new resource is effectively frozen out forever.

\emph{2. Content-based prior.} Encode the new quiz's metadata---topic (``linear algebra''), difficulty level, length---and find similar resources that already have engagement. Use the average click-through rate (CTR) of those neighbors as a prior CTR for the new quiz. This warmstarts the prediction. The bet is that the metadata is informative and that the new resource will behave like its neighbors.

\emph{3. Exploration bonus.} Add an explicit bonus $b/\sqrt{n_i + 1}$ to every resource's predicted score, where $n_i$ is the number of times resource $i$ has been recommended so far. This guarantees the new quiz gets surfaced until it accumulates enough data to be evaluated on its own merit. The bet is short-term loss of CTR (sometimes showing items that may not be top-rated) in exchange for long-term catalog coverage.

\emph{Concrete numbers.} Suppose popularity fallback gives the current top three resources predicted CTRs of $0.18$, $0.16$, $0.15$. The new quiz has no observed data. Content-based assigns it a prior CTR of $0.14$ from the average of similar quizzes. With exploration bonus $b = 0.05$ and $n_i = 0$, the bonus adds $0.05 / \sqrt{1} = 0.05$. The new quiz's score becomes $0.14 + 0.05 = 0.19$, just above the popularity leader at $0.18$, so it gets recommended. After $20$ impressions and $4$ clicks, the bonus has shrunk to $0.05 / \sqrt{21} \approx 0.011$, and the resource is now evaluated essentially on its observed CTR of $4/20 = 0.20$.

Each strategy makes a different bet about what should replace the missing cold-start data. The popularity strategy bets that new items will never matter. The content-based strategy bets on metadata quality. The exploration-bonus strategy bets that paying short-term CTR for long-term coverage is the right trade.
\end{example}

\section{Matrix Factorization}

\emph{Back to the course-support recommender.} User- and item-based collaborative filtering gave the team a working baseline but two persistent problems: cold-start students with no interaction history get nothing useful, and similarity over a sparse 500-resource catalog is noisy. Matrix factorization is what the team reaches for next --- a way to learn compact ``what this student tends to need'' and ``what this resource tends to teach'' vectors from the same logs, with explicit handles for regularization and cold-start fallbacks.
\label{sec:matrix-factorization}

\subsection{The Latent Factor Model}

Start with the concrete object being modeled: a sparse user-item table. A blank entry is not a rating of zero; it usually means the system did not observe that user interacting with that item.

\begin{center}
\begin{tabular}{lcccc}
\toprule
 & Resource A & Resource B & Resource C & Resource D \\
\midrule
Student 1 & 5 &  & 4 &  \\
Student 2 &  & 4 &  & 1 \\
Student 3 & 4 &  &  & 5 \\
\bottomrule
\end{tabular}
\end{center}

The pedagogical question is whether the observed pattern can suggest a plausible score for a missing entry, such as Student 1 on Resource B, without pretending that every blank is negative feedback. Matrix factorization answers by giving each student and each resource a small latent vector, then using vector compatibility to fill in plausible scores for unobserved pairs.

Matrix factorization addresses the sparsity problem by learning latent factors for users and items. The interaction matrix $R$ (users $\times$ items) is approximated as the product of two low-rank matrices:
\[
R \approx U V^\top
\]

where $U \in \mathbb{R}^{m \times k}$ holds $k$-dimensional latent vectors for each of $m$ users, and $V \in \mathbb{R}^{n \times k}$ holds $k$-dimensional latent vectors for each of $n$ items. The predicted rating for user $u$ on item $i$ is:
\[
\hat{r}_{ui} = u_u^\top v_i
\]

The latent factors are learned by minimizing the reconstruction error on observed entries:
\[
\min_{U, V} \sum_{(u,i) \in \text{observed}} (r_{ui} - u_u^\top v_i)^2 + \lambda(\|U\|_F^2 + \|V\|_F^2)
\]

The regularization term $\lambda(\|U\|_F^2 + \|V\|_F^2)$ uses squared Frobenius norms---the sums of squared entries in the user and item factor matrices. It prevents overfitting; without it, the model can memorize the observed ratings without learning generalizable structure.

\textbf{Two ways to optimize the factorization.} The objective is non-convex jointly in $U$ and $V$ but \emph{convex} in each one with the other held fixed. That observation is the basis of both standard optimizers.

\emph{Alternating least squares (ALS).} Fix $V$ and minimize over $U$: each user row $u_u$ has a closed-form ridge-regression update obtained by setting the gradient to zero. Writing $\Omega_u = \{i : (u, i)\text{ observed}\}$,
\[
u_u \;\leftarrow\; \Bigl(\sum_{i \in \Omega_u} v_i v_i^\top + \lambda I\Bigr)^{-1} \sum_{i \in \Omega_u} r_{ui}\, v_i.
\]
Then fix $U$ and apply the symmetric update for each item $v_i$. Alternate to convergence. Each half-step is a stack of $k \times k$ linear systems, one per user or item, and is trivially parallel.

\emph{Stochastic gradient descent (SGD).} For one observed $(u,i)$, write the prediction error $\epsilon_{ui} = r_{ui} - u_u^\top v_i$. The gradient of the loss on this single example gives the paired updates
\[
u_u \;\leftarrow\; u_u + \eta\,(\epsilon_{ui}\, v_i - \lambda\, u_u),
\qquad
v_i \;\leftarrow\; v_i + \eta\,(\epsilon_{ui}\, u_u - \lambda\, v_i),
\]
applied for one sampled observation at a time. SGD is preferred when the observed matrix is too large to materialize the per-user systems, and when implicit-feedback weights or side features make the closed form intractable.

The two optimizers reach essentially the same fixed points on the same objective; ALS converges in fewer outer iterations but pays more per iteration, while SGD scales smoothly to streaming data.

\begin{codeexample}
import numpy as np

def als_step(R, mask, U, V, lam):
    """One full ALS pass over a sparse rating matrix R (m, n), with
    a boolean mask of observed entries (m, n)."""
    m, n = R.shape
    k    = U.shape[1]
    Ik   = lam * np.eye(k)
    # Update U row-by-row given V
    for u in range(m):
        idx  = np.where(mask[u])[0]
        if len(idx) == 0:
            continue
        Vi   = V[idx]                       # (|Omega_u|, k)
        A    = Vi.T @ Vi + Ik               # (k, k)
        b    = Vi.T @ R[u, idx]             # (k,)
        U[u] = np.linalg.solve(A, b)
    # Update V symmetrically given U
    for i in range(n):
        idx  = np.where(mask[:, i])[0]
        if len(idx) == 0:
            continue
        Uu   = U[idx]                       # (|Omega_i|, k)
        A    = Uu.T @ Uu + Ik
        b    = Uu.T @ R[idx, i]
        V[i] = np.linalg.solve(A, b)
    return U, V

# 4 users, 5 items, partially observed
rng  = np.random.RandomState(0)
R    = rng.uniform(1, 5, size=(4, 5))
mask = rng.uniform(size=(4, 5)) < 0.6        # ~60% observed
R   *= mask                                  # zero out unobserved
k, lam = 2, 0.1
U = rng.normal(scale=0.1, size=(4, k))
V = rng.normal(scale=0.1, size=(5, k))
for _ in range(20):
    U, V = als_step(R, mask, U, V, lam)
print((U @ V.T).round(2))                    # filled-in predictions
\end{codeexample}

The structure is the algorithm: hold one factor matrix fixed, solve a per-user (or per-item) ridge regression on its observed entries, then swap and repeat. Confidence-weighted ALS for implicit feedback simply replaces the unweighted $V_i^\top V_i$ with $V_i^\top C_i V_i$, where $C_i$ is a diagonal confidence matrix derived from interaction strength.

\subsection{What the Latent Factors Mean}

The latent factors capture implicit structure. In a movie recommendation system, one factor might encode ``action vs.\ drama preference,'' another ``preference for recent vs.\ classic films.'' In a course-resource system, one factor might encode ``prefers worked examples vs.\ conceptual explanations,'' another ``prefers short vs.\ long content.'' The factors are not predefined---they are discovered from the interaction data.

\begin{example}[Matrix Factorization for Course Resources]
The course-support system has a sparse interaction matrix: 800 students $\times$ 500 resources, with only 16,000 observed interactions (4\% fill rate). Matrix factorization with $k=20$ latent factors learns user vectors that capture resource-format and study-pattern preferences, and item vectors that capture resource characteristics. After training, the system predicts that Student~A---whose latent vector resembles students who prefer worked examples---would benefit from Resource~247, a worked example on regression, even though Student~A has never accessed it. The predicted relevance score is 4.2 out of 5, placing it in the student's top-5 recommendations.
\end{example}

Matrix factorization is conceptually related to the dimensionality reduction in Chapter~\ref{ch:structure-beyond-linearity}. PCA also factorizes a matrix into low-rank components. The key difference is that matrix factorization for recommendation works on a \emph{sparse} matrix (most entries are missing) and fits only the observed entries, while PCA operates on a complete matrix. The optimization algorithms differ accordingly: alternating least squares (ALS) or stochastic gradient descent (SGD) for sparse matrix factorization, versus eigendecomposition for PCA.

\subsection{Adding Side Information}

Pure matrix factorization uses only the interaction data, but additional information is often available: user demographics, item descriptions, temporal patterns. \emph{Hybrid models} combine collaborative filtering with content-based features:

\begin{itemize}[leftmargin=*]
\item \emph{Content-based features.} Use item attributes (resource type, topic, difficulty level) to compute item similarity, addressing the cold-start problem for new items.
\item \emph{Context features.} Use temporal and contextual signals (time of day, device, position in the semester) to adjust recommendations.
\item \emph{Deep learning models.} Neural collaborative filtering replaces the dot product $u_u^\top v_i$ with a neural network that takes user and item embeddings as input. This can capture nonlinear interactions but requires more data and is harder to interpret.
\end{itemize}

\section{Implicit Feedback: The Missing-Data Problem}
\label{sec:implicit-feedback}

\subsection{Explicit vs.\ Implicit Signals}

Most real-world recommendation systems lack explicit ratings (1--5 stars). Instead, they observe \emph{implicit feedback}: clicks, views, dwell time, downloads, and purchases. Implicit feedback is abundant but ambiguous.

\textbf{The asymmetry problem.} A click signals some interest, but a non-click is ambiguous. A student who did not access a resource may have been uninterested---or may never have seen it. A user who did not click a movie may dislike it---or may not have scrolled far enough. Implicit feedback provides positive signals (``the user interacted with this'') but no reliable negative signals (``the user saw this and chose not to interact'').

\textbf{The noise problem.} Implicit signals are noisy. A student may click a resource by accident. A user may watch a movie for 30 seconds and abandon it. Dwell time on a webpage may signal engagement---or confusion. Turning implicit signals into training labels requires careful design.

\subsection{Treating Implicit Feedback as Training Data}

Common strategies for learning from implicit feedback:

\begin{itemize}[leftmargin=*]
\item \emph{Binary encoding.} Treat all interactions as positive (1) and all non-interactions as negative (0). Simple, but it treats ambiguous non-interactions as definitive negatives.
\item \emph{Confidence weighting.} Assign higher confidence to observed interactions and lower confidence to non-interactions. The model treats observed interactions as stronger evidence of preference and unobserved ones as low-confidence zero-preference entries rather than certain dislikes. This is the approach used in the influential ``implicit ALS'' framework.
\item \emph{Negative sampling.} Instead of treating all non-interactions as negatives, randomly sample a subset as negatives. This cuts computational cost and avoids overwhelming the model with uninformative negatives.
\item \emph{Exposure modeling.} Explicitly model whether the user \emph{saw} the item. If the item was never displayed, the non-interaction is missing data, not a negative signal. This requires logging what was recommended, which connects to the exposure bias discussion below.
\end{itemize}

\begin{warningbox}
Treating all non-interactions as negatives is one of the most common mistakes in recommendation system design. If the system has shown a user only 50 of 500 available items, the 450 unseen items are unknown, not negative. Training on them as negatives teaches the model to suppress items the user has never had a chance to evaluate. This systematically biases the model toward re-recommending items similar to what has already been shown.
\end{warningbox}

\begin{example}[Implicit Feedback in Course Support]
The course-support system logs which resources each student accesses, how long they spend on each, and whether they return. A student who spends 45 minutes on a worked example and returns to it twice is a strong positive signal. A student who clicks a resource and leaves after 10 seconds is a weak positive (or possibly a negative). A resource the student was never shown is missing data, not a negative. The team uses confidence-weighted matrix factorization: observed interactions are weighted by dwell time (longer means higher confidence), and non-interactions for items that were \emph{displayed but not clicked} are treated as weak negatives. Non-interactions for items that were never displayed are excluded from training.
\end{example}

\subsection{Listwise Losses: When Pairs Aren't Enough}

Matrix factorization with squared error is a \emph{pointwise} loss: each $(u, i)$ prediction is scored independently against its observed rating. Bayesian Personalized Ranking (BPR) and RankNet are \emph{pairwise} losses: each pair $(i^+, i^-)$ of a relevant and a non-relevant item for the same user is treated as a binary classification problem---``did the model rank $i^+$ above $i^-$?'' Both losses are convenient, but neither captures the operational truth that ranking is about the \emph{full ordered list} the user actually sees, and that an error at position $1$ matters far more than an error at position $20$.

\emph{Listwise approaches} put the whole list into the loss. LambdaRank (Burges, 2010) keeps the pairwise scaffolding but weights each pair's gradient by its impact on a position-sensitive metric: if swapping the two items would change NDCG by $\Delta\text{NDCG}$, the pair's gradient is scaled by $|\Delta\text{NDCG}|$. Pairs whose swap would barely move the metric get little weight; pairs whose swap would move a high-relevance item out of the top positions get the most. ListNet (Cao et al., 2007) goes further: it treats the full ranking as a probability distribution over permutations using a Plackett-Luce parameterization, and minimizes the KL-divergence between the model's induced permutation distribution and the ideal one. Both are operational refinements of the same statement: \emph{I care about the ranked list, not the pair.}

\emph{Choosing among the three.} Use pointwise when ratings are explicit and calibration matters---predicting an absolute score that another downstream system will consume. Use pairwise (BPR) when you have implicit clicks or conversions and want a simple binary preference signal that scales easily with negative sampling. Use listwise (LambdaRank, ListNet) when you have an exposure-aware top-$k$ metric---NDCG, MRR, Recall@$k$---that you actually report to stakeholders; the position-weighting in the loss is what aligns training with reporting.

The choice of pointwise / pairwise / listwise is a modeling choice about \emph{what counts as the same answer}. Pointwise loss says two rankings are the same when their scores match. Pairwise loss says they are the same when their pairwise orderings match. Listwise loss says they are the same when their top-$k$ items appear in the same positions. Each definition aligns with a different reporting practice, and the right choice is the one whose ``same'' matches what the user experiences.

\section{Evaluating Ranking Systems}
\label{sec:ranking-eval}

\subsection{Position-Aware Metrics}

Classification metrics (accuracy, precision, recall) miss what matters in ranking: the \emph{order} of the results. A recommendation list that places the best item at position 10 is worse than one that places it at position 1, even when both contain the same items.

\begin{example}[A Tiny Ranked-List Walkthrough]
Suppose the course-support system shows a student three resources in this order: \#1 worked example, \#2 short video, \#3 practice quiz. Later we learn the worked example was somewhat helpful, the video was not helpful, and the quiz was very helpful. With binary relevance, the relevant items sit at ranks 1 and 3. The reciprocal rank is $1$ because the first relevant item is at rank 1, so MRR for this list is perfect even though another relevant item is lower. Average Precision looks at both relevant positions: precision@1 is $1/1=1$ and precision@3 is $2/3$, and AP averages them. With graded relevance instead---say scores 2, 0, and 3---NDCG gives more credit when the very helpful quiz is near the top and less when it is buried lower in the list.
\end{example}

For the same list, binary AP is
\[
\mathrm{AP}=\frac{1+\frac{2}{3}}{2}\approx 0.833.
\]
With graded relevance \(2,0,3\), the discounted gain is
\[
\mathrm{DCG@3}=2+\frac{0}{\log_2 3}+\frac{3}{\log_2 4}=3.5.
\]
The ideal order places relevance \(3\) first and \(2\) second, giving
\[
\mathrm{IDCG@3}=3+\frac{2}{\log_2 3}\approx 4.262,
\qquad
\mathrm{NDCG@3}\approx 0.821.
\]
The numbers make the differences visible: MRR is perfect because the first relevant item appears first, AP notices that another relevant item is lower, and NDCG notices that the most helpful item was not first.

\textbf{Mean Reciprocal Rank (MRR).} The reciprocal rank of a single query is $1/\text{rank}$ of the first relevant item. MRR averages this across queries. MRR rewards systems that place the first relevant item high but ignores all other relevant items.

\textbf{Mean Average Precision (MAP).} Average Precision for a single query computes precision at each position where a relevant item appears, then averages those values. MAP averages across queries. It rewards systems that place \emph{all} relevant items high, not just the first.

\textbf{Normalized Discounted Cumulative Gain (NDCG).} NDCG assigns graded relevance scores to items rather than binary relevant/irrelevant labels and discounts items lower in the list. Writing $\mathrm{rel}_i$ for the graded relevance of the item at rank $i$, the discounted cumulative gain at depth $k$ is
\[
\mathrm{DCG}@k \;=\; \sum_{i=1}^{k} \frac{2^{\mathrm{rel}_i} - 1}{\log_2(i + 1)},
\]
and the normalized version divides by the ideal value $\mathrm{IDCG}@k$, the DCG of the same items reordered by descending relevance:
\[
\mathrm{NDCG}@k \;=\; \frac{\mathrm{DCG}@k}{\mathrm{IDCG}@k} \;\in\; [0, 1].
\]
The numerator $2^{\mathrm{rel}_i} - 1$ is the exponential gain convention used in industry; the logarithmic position discount $1/\log_2(i + 1)$ is the standard rank discount. The normalization makes NDCG comparable across queries with different numbers of relevant items. NDCG is the most widely used ranking metric because it handles graded relevance and position discounting together.

NDCG@$k$ evaluates only the top $k$ positions. Choosing $k$ is a design decision: for a mobile app showing 5 recommendations, NDCG@5 is appropriate; for a search page showing 20 results, NDCG@20 may be better. Always report $k$ and justify it from the user interface.

\subsection*{Extension: Counterfactual Offline Evaluation (IPS and Doubly Robust)}

This subsection introduces inverse-propensity-scoring (IPS) and doubly-robust (DR) estimators for offline evaluation. The estimators are standard in modern recommendation but use the expectation-under-a-distribution language from Mathematical Bridge~II and an importance-weighting argument that takes a section's worth of setup. The headline is short: naive offline replay is biased by the old exposure policy, and principled estimators reweight the logs to correct for that. The Decision Box at the end of the section captures what the chapter artifact actually needs; readers who want only that operational summary can skip the derivations and read the formulas as recipes.

In logged implicit-feedback data, naive offline evaluation is unreliable when the test observations were generated by a previous exposure policy. Imagine last semester's system almost always showed worked examples first and rarely showed quizzes. If a student clicked the worked example, the log records that interaction. But if the student would have preferred the quiz, there is no evidence of that because the quiz was never exposed. This is the \emph{missing-not-at-random} problem: the test data is biased by the policy that generated it.

\textbf{Counterfactual estimation.} The most principled approach uses inverse propensity scoring (IPS). Let $\pi_0$ be the old logging policy that generated the data and $\pi$ the new target policy whose value we want to estimate. For each logged interaction $(u, i)$ with observed outcome $y_{ui}$ (a click, a reward, a relevance label), the IPS estimator of the new policy's value is
\[
\hat{V}_{\text{IPS}}(\pi) \;=\; \frac{1}{N}\sum_{(u, i)\in\mathcal{D}} \frac{\pi(i \mid u)}{\pi_0(i \mid u)}\, y_{ui},
\]
where $\pi_0(i \mid u)$ is the logged propensity---the probability that the old system would have exposed item $i$ to user $u$. The intuition is that a logged interaction tells us about whatever $\pi_0$ chose to expose, and dividing by $\pi_0$ rescales each datum to be representative under the new policy. If quizzes were shown only 10\% of the time, a click on a quiz carries weight $1/0.1 = 10$; a click on a worked example shown 90\% of the time carries weight $1/0.9 \approx 1.1$. Items the previous system was unlikely to show get higher weight, correcting the selection bias. Standard estimator properties: $\hat{V}_{\text{IPS}}$ is unbiased when $\pi$ has support inside $\pi_0$, but its variance grows with $1/\pi_0$ on the rare side, which is why small propensities are typically clipped at a floor $\pi_0 \ge \pi_{\min}$ (a bias--variance trade with a small introduced bias). This requires logging item-position exposure propensities, which many production systems do not do.

\begin{example}[$\hat V_{\text{IPS}}$ on a Five-Row Log]
A logged interaction set, with reward $r_i = 1$ if the user clicked and $0$ otherwise, and logged propensity $\pi_0$:
\begin{center}
\begin{tabular}{cllcc}
\toprule
$i$ & user & item & $r_i$ & $\pi_0(a_i\mid x_i)$ \\
\midrule
$1$ & $u_1$ & quiz A   & $1$ & $0.10$ \\
$2$ & $u_2$ & quiz B   & $0$ & $0.50$ \\
$3$ & $u_3$ & video C  & $1$ & $0.20$ \\
$4$ & $u_4$ & quiz A   & $1$ & $0.40$ \\
$5$ & $u_5$ & video D  & $0$ & $0.05$ \\
\bottomrule
\end{tabular}
\end{center}
Define the new target policy as $\pi(a_i\mid x_i) = \min(2\,\pi_0(a_i\mid x_i),\,1)$, so the new system is twice as likely to recommend each logged item, capped at $1$. The per-row importance weight is $w_i = \pi(a_i\mid x_i)/\pi_0(a_i\mid x_i)$ (which equals $2$ except when capping bites), and the contribution to the estimator is $w_i r_i$.
\begin{center}
\begin{tabular}{ccccc}
\toprule
$i$ & $\pi_0$ & $\pi$ & $w_i = \pi/\pi_0$ & $w_i\, r_i$ \\
\midrule
$1$ & $0.10$ & $0.20$ & $2.0$ & $2.0$ \\
$2$ & $0.50$ & $1.00$ & $2.0$ & $0.0$ \\
$3$ & $0.20$ & $0.40$ & $2.0$ & $2.0$ \\
$4$ & $0.40$ & $0.80$ & $2.0$ & $2.0$ \\
$5$ & $0.05$ & $0.10$ & $2.0$ & $0.0$ \\
\bottomrule
\end{tabular}
\end{center}
Summing and averaging over $n=5$ gives $\hat V_{\text{IPS}} = (2.0 + 0 + 2.0 + 2.0 + 0)/5 = 1.2$. To see why variance is the real concern, swap row $1$'s reward into a counterfactual log where the click had landed on the rarely-exposed video D ($\pi_0 = 0.05$): that single row would contribute $w_5 r_5 = (0.10/0.05)\cdot 1 = 2$, fine here, but if the target policy had assigned $\pi = 0.5$ to that pair the weight jumps to $10$ and a single click moves the estimator by $10/5 = 2$ in absolute value. Clipping weights at, say, $w_i \le 5$ trades a small bias for a large variance reduction, which is why production IPS deployments almost always clip.
\end{example}

\begin{decisionbox}
IPS is only as credible as the logging policy behind it. Before using it, verify three things. First, the old system logged which items were exposed, their positions or slate context, and the probability of each exposure. Second, the new policy mostly recommends items inside the old policy's support. Third, very small propensities are clipped, stabilized, or analyzed for variance. Without these conditions, IPS can look principled while producing a noisy or unsupported estimate.
\end{decisionbox}

\textbf{Doubly-robust estimation.} A practical refinement combines IPS with a learned reward model $\hat r(u, i)$ that predicts the outcome from features. The doubly-robust estimator subtracts the model prediction from the importance-weighted residual and adds back the model's expected value under the new policy:
\[
\hat{V}_{\mathrm{DR}}(\pi) \;=\; \frac{1}{N}\sum_{(u, i) \in \mathcal{D}} \Bigl[\, \hat{r}_\pi(u) \;+\; \frac{\pi(i \mid u)}{\pi_0(i \mid u)}\bigl(y_{ui} - \hat r(u, i)\bigr)\Bigr],
\qquad \hat r_\pi(u) = \sum_{i'} \pi(i' \mid u)\, \hat r(u, i').
\]
The name describes the guarantee: $\hat{V}_{\mathrm{DR}}$ is unbiased if \emph{either} the propensities $\pi_0$ are correct \emph{or} the reward model $\hat r$ is correct, not necessarily both. When the reward model is accurate, the residual $y_{ui} - \hat r(u, i)$ is small, the importance-weighted term contributes little variance, and the estimator reduces to the model's predictions $\hat r_\pi(u)$---this is the variance-reduction benefit over plain IPS. When the propensities are correct and the model is wrong, the residual term still corrects the model's bias on the data the new policy would actually expose. The DR estimator is the standard counterfactual evaluator in modern recommendation work because it tolerates one source of misspecification at a time, which is realistic in deployed systems where neither propensities nor reward models are perfectly known.

\textbf{Replay methods.} When historical logs include randomized recommendations---say, a small fraction of random items mixed into the ranked list---replay methods estimate how a new policy would have performed by replaying the historical data and using the randomized portion as unbiased logged observations. In the same example, if the old system occasionally inserted a random quiz into the list, those randomized exposures provide a fairer basis for judging whether a new policy that ranks quizzes higher would help.

\begin{warningbox}
Offline metrics on logged implicit-feedback data tend to favor policies similar to the logger and understate policies that recommend items outside its support. Drawing strong conclusions requires randomized exposure, explicit relevance judgments, or counterfactual correction. Complement offline evaluation with online A/B tests before concluding that a new system is worse.
\end{warningbox}

\section{Exposure Bias and Feedback Loops}
\label{sec:exposure-bias}

\emph{The dilemma from the opener, named.} The course-support recommender that ``taught itself to recommend only what it already showed'' is not a freak failure mode. It is the natural behaviour of any system that trains on logs of its own past decisions. This section names the mechanism, says why the IPS and DR estimators of the previous section exist, and gives the audit habits the chapter artifact will codify.

\subsection{How Recommender Systems Shape Their Own Data}

A classifier passively receives test inputs. A recommender system actively decides what users see. That difference creates a feedback loop:

\begin{enumerate}[leftmargin=*]
\item The system recommends certain items to users.
\item Users interact with the recommended items (because those are what they see).
\item The interactions become the training data for the next version of the system.
\item The next version learns to recommend items similar to what was previously recommended.
\end{enumerate}

Over time, the system converges on a narrow slice of the item space. Popular items get more exposure, collect more data, and become even more popular. New, niche, or minority-interest items receive less exposure, collect less data, and grow invisible.

\subsection{Consequences of Exposure Bias}

\textbf{Popularity amplification.} A few items dominate recommendations regardless of individual preferences. The ``long tail'' of less popular but potentially relevant items is suppressed.

\textbf{Filter bubbles.} Users are shown increasingly narrow content that reinforces past behavior. A student who once accessed a worked example on regression is shown more regression examples, even after they have moved on.

\textbf{Fairness to providers.} On two-sided platforms (job boards, course marketplaces, content platforms), items are created by providers---instructors, content creators, job posters. Exposure bias can systematically disadvantage providers whose items were not recommended early, creating unfair competitive dynamics.

\textbf{Evaluation corruption.} The feedback loop corrupts the evaluation data. The system looks good in offline evaluation because it is tested on data it generated. In absolute terms it may be performing poorly: users are satisfied only because they were never shown better alternatives.

\begin{example}[Exposure Bias in Campus Job Recommendations]
A university job board recommends positions to graduating students based on past click data. Early on, software engineering jobs received the most clicks because they were listed first alphabetically. The system learned that software engineering is ``popular'' and began recommending it more. Over time, research assistant positions, nonprofit roles, and teaching opportunities received less exposure, fewer clicks, and even less recommendation. A student interested in teaching would have to search actively rather than discover opportunities through recommendations. The feedback loop turned the system into a funnel toward software engineering, regardless of individual interests.
\end{example}

\subsection{Mitigating Exposure Bias}

\begin{itemize}[leftmargin=*]
\item \emph{Exploration.} Deliberately recommend some items the system is uncertain about, not just items it is confident the user will like. This mirrors the exploration-exploitation tradeoff in reinforcement learning (Chapter~\ref{ch:reinforcementlearning}).
\item \emph{Diversity constraints.} Require that the top-$k$ items span multiple categories, topics, or providers. This prevents the list from converging on a single niche.
\item \emph{Exposure fairness.} Define a target exposure policy across eligible providers or content categories, then monitor deviations from that target while accounting for relevance, inventory, and user needs. This connects to the fairness discussion in Chapter~\ref{ch:reliability}.
\item \emph{Logging and counterfactual evaluation.} Log what was recommended and with what probability. Use inverse propensity scoring to evaluate alternative policies on historical data. This catches exposure bias before it becomes entrenched.
\end{itemize}

\begin{decisionbox}
\textbf{Exposure bias checklist:}
\begin{itemize}
\item Does the system log what it recommends and with what probability?
\item Is there an exploration mechanism (randomization, $\epsilon$-greedy, Thompson sampling)?
\item Are diversity or fairness constraints enforced on the recommendation list?
\item Has the system been evaluated with counterfactual methods, not just offline replay?
\item Is the feedback loop monitored: is the item distribution in training data narrowing over time?
\end{itemize}
\end{decisionbox}

\subsection{Bandits and Contextual Bandits}

When an exploration strategy is built into the recommendation system itself rather than added as a one-off randomization, the natural framework is the \emph{multi-armed bandit}. A bandit treats each candidate item (or arm) as offering a stochastic reward; the policy must balance pulling arms that already look good against pulling arms whose value is still uncertain. Algorithms include $\epsilon$-greedy, Upper Confidence Bound (UCB) \citep{auer2002ucb}, and Thompson sampling \citep{chapelle2011thompson}. \emph{Contextual bandits} extend this by conditioning the arm choice on observed user features, which makes them the natural minimal model of online recommendation: each request is a context, each candidate is an arm, and the system learns from the click or non-click that follows. Chapter~\ref{ch:reinforcementlearning} introduced this tradeoff in the full RL setting; for many production recommendation problems the contextual-bandit subset is enough.

\section{Chapter Artifact: The Recommendation Audit Card}
\label{sec:rec-audit}

The chapter artifact is the \emph{recommendation audit card}: a structured document that records what the system recommends, which signals it uses, what it exposes and suppresses, and which feedback loops exist. It connects to the data design card (Chapter~\ref{ch:dataset-design}), the uncertainty audit card (Chapter~\ref{ch:probabilistic-models}), and the deployment considerations in Part~V.

\subsection{Recommendation Audit Card Structure}

\begin{table}[t]
\centering
\small
\setlength{\tabcolsep}{4pt}
\begin{tabular}{@{}L{2.8cm}L{7.2cm}@{}}
\toprule
Card field & What must be documented \\
\midrule
Recommendation task & What is being ranked, for whom, and in what context \\
Signal sources & What user signals drive ranking (explicit ratings, clicks, dwell time, purchases, etc.) \\
Model family & Collaborative filtering, matrix factorization, content-based, hybrid, or neural \\
Cold-start strategy & How new users and new items are handled (defaults, content features, exploration) \\
Implicit feedback treatment & How non-interactions are treated (binary, confidence-weighted, exposure-modeled) \\
Evaluation protocol & What metrics are used (NDCG@$k$, MAP, MRR), offline vs.\ online evaluation, counterfactual correction \\
Exposure analysis & What items or categories are overrepresented in recommendations; what is suppressed \\
Diversity and fairness & Whether diversity or fairness constraints are enforced; how provider-side fairness is measured \\
Feedback loop risk & Whether the system generates its own training data; what monitoring is in place \\
Exploration mechanism & Whether and how the system explores uncertain or underrepresented items \\
\bottomrule
\end{tabular}
\caption{A recommendation audit card documents what the system exposes, suppresses, and amplifies.}
\label{tab:rec-audit-card}
\end{table}

\subsection{Filled Example: Course Resource Recommender}

\begin{example}[Recommendation Audit Card for Course Resource System]

\textbf{Recommendation task.} Rank 500 learning resources for students flagged as CONCERNING or URGENT by the early-warning model. Top-5 recommendations appear on the student dashboard.

\textbf{Signal sources.} Implicit feedback only: resource access logs (binary click), dwell time (seconds), and return visits (count). No explicit ratings are collected.

\textbf{Model family.} Hybrid: confidence-weighted matrix factorization ($k=20$ latent factors) combined with content-based features (resource type, topic tag, difficulty level). Content features address cold start for newly uploaded resources.

\textbf{Cold-start strategy.} New students receive recommendations based on content features and the course's most-accessed resources (popularity fallback). New resources are shown to a random 5\% of users for the first two weeks to gather initial interaction data (exploration).

\textbf{Implicit feedback treatment.} Confidence-weighted: access events are weighted by $1 + \alpha \cdot \log(1 + t/t_0)$, where $t$ is dwell time in seconds, $t_0=60$ seconds is the reference duration, and $\alpha = 0.5$. Items displayed but not clicked are weak negatives (confidence 0.1). Items never displayed are excluded from training.

\textbf{Evaluation protocol.} NDCG@5 and MAP@5 on a held-out set of student-resource interactions from the most recent semester. Counterfactual correction uses IPS with logged item-position exposure propensities for the displayed slate. An online A/B test is planned for fall semester.

\textbf{Exposure analysis.} Worked examples and lecture recordings account for 68\% of all recommendations despite making up only 40\% of the library. Office-hour transcripts and tutoring sessions are underrepresented. Resources uploaded in the current semester receive 3$\times$ less exposure than older resources because they have less interaction data.

\textbf{Diversity and fairness.} The top-5 list must include at least 2 distinct resource types (it cannot show 5 worked examples, for example). No instructor-side fairness constraint is currently enforced; this is a known gap.

\textbf{Feedback loop risk.} High. The system generates the data it trains on. Worked examples dominate because they were recommended early, collected more clicks, and are now recommended even more. Monitoring: track normalized entropy of the resource-type distribution in training data quarterly, where 1 means perfectly balanced exposure across resource types and 0 means all exposure goes to one type. Alert if normalized entropy drops below 0.75, indicating dangerous narrowing.

\textbf{Exploration mechanism.} $\epsilon$-greedy with $\epsilon = 0.05$: 5\% of recommendation slots are filled with a random resource the student has not yet seen. New resources receive a 2-week boost, shown to 5\% of users regardless of score.

\end{example}

\section{Common Mistakes in Recommendation Systems}

Several mistakes recur when building and deploying recommender systems. The first is treating all non-interactions as negative signals, which biases the model toward re-recommending what it has already shown. The second is evaluating offline without counterfactual correction, which makes the current system look artificially good. The third is ignoring the cold-start problem: if new items are never recommended, they can never collect the data needed to be recommended, creating a permanent invisibility trap. The fourth is neglecting exposure fairness: a system that maximizes click-through rate may systematically suppress content from minority providers, niche topics, or newly created resources. Finally, deploying a recommender without monitoring the feedback loop---checking whether the item distribution in training data is narrowing over time---risks a system that converges on a narrow, self-reinforcing slice of the item space.

\section{Looking Ahead}

This chapter treated recommendation as a ranking problem with distinctive challenges: implicit feedback, missing data, exposure bias, and feedback loops. The recommendation audit card documents what the system exposes and suppresses. These challenges preview the broader themes of Part~V: experiments that produce honest claims (Chapter~\ref{ch:experiments-claims}), reliability under distribution shift (Chapter~\ref{ch:reliability}), and the transition from models to production systems whose outputs shape the world they operate in (Chapter~\ref{ch:models-systems}).

\section*{Chapter Summary}

Recommendation is a ranking problem: the system produces an ordered list, and the order determines what users see. Collaborative filtering recommends items liked by similar users but fails at cold start. Matrix factorization learns latent factors for users and items, addressing sparsity at the cost of requiring regularization. Implicit feedback (clicks, dwell time) is abundant but ambiguous: non-interactions are missing data, not negative signals. Position-aware metrics (NDCG, MAP, MRR) evaluate ranking quality. Offline evaluation is biased by the previous system's recommendations and requires counterfactual correction. Exposure bias creates feedback loops: the system shapes its own training data, amplifying popular items and suppressing the long tail. Mitigation requires deliberate exploration, diversity constraints, and feedback-loop monitoring. The recommendation audit card documents what the system exposes, suppresses, and amplifies.

\begin{studyquestions}
\item How does recommendation differ from classification, and why does the difference matter?
\item What is the cold-start problem, and why is it a design problem rather than just a technical one?
\item What do the latent factors in matrix factorization represent?
\item Why are non-interactions in implicit feedback ambiguous, and what are the consequences of treating them as negatives?
\item Why does NDCG discount items lower in the ranked list?
\item Why can offline evaluation on logged implicit-feedback data be biased, and what logging or randomization reduces the problem?
\item How does a recommender system create a feedback loop, and what are the consequences?
\item What is exposure bias, and how does it differ from popularity bias?
\item Why is exploration important in recommendation, and how does it connect to reinforcement learning?
\item What should a recommendation audit card document about feedback loops?
\end{studyquestions}

\begin{exercises}
\item \exercisetag{Calculation}{Core}Compute NDCG@5 by hand for two ranked lists that contain the same 5 items in different orders. Show how position affects the score.
\item \exercisetag{Design}{Core}For the tiny user-item table in this section, explain why a blank entry should not automatically be treated as dislike. Name one measurement that would distinguish non-exposure from exposed-but-ignored.
\item \exercisetag{Design}{Intermediate}Design an implicit feedback scheme for a course resource recommender. Specify: what constitutes a positive signal, how non-interactions are handled, and what confidence weights are used. Justify each decision.
\item \exercisetag{Design}{Intermediate}Design an exploration strategy for a job board recommender that must balance relevance (showing jobs the student is likely to want) with fairness (ensuring new job postings get enough exposure to collect initial click data). Specify the tradeoff parameter and how you would tune it.
\item \exercisetag{Design}{Intermediate}Construct a recommendation audit card (Table~\ref{tab:rec-audit-card}) for a recommendation system you use daily (e.g., a streaming service, a social media feed, a shopping site). Identify at least two feedback loop risks.
\item \exercisetag{Implementation}{Advanced}Given a sparse user-item interaction matrix (e.g., MovieLens-100K), implement matrix factorization with SGD. Vary the number of latent factors $k \in \{5, 10, 20, 50\}$ and plot RMSE on a held-out test set. At what $k$ does overfitting begin?
\item \exercisetag{Implementation}{Advanced}For the same dataset, compare user-based collaborative filtering (cosine similarity) with matrix factorization on NDCG@10. Which performs better, and why?
\end{exercises}

\begin{challengeexercises}
\item \exercisetag{Diagnosis}{Advanced}A course resource recommender shows worked examples 3$\times$ more than other resource types. A team member argues this is because students genuinely prefer worked examples. A critic argues this is because the system recommended worked examples early and the feedback loop amplified them. Design an experiment to distinguish between these explanations.
\item \exercisetag{Diagnosis}{Advanced}Explain why inverse propensity scoring corrects for selection bias in offline evaluation. Under what conditions does IPS have high variance, and what can be done about it?
\item \exercisetag{Diagnosis}{Advanced}A university deploys a job recommender and finds that after 6 months, 90\% of student clicks go to 10\% of listed jobs. Diagnose whether this is a preference concentration problem (students genuinely prefer a few job types) or an exposure bias problem (the system only showed a few job types). Propose three specific measurements.
\item \exercisetag{Design}{Advanced}Design a recommender system for a two-sided marketplace (e.g., tutors and students) that optimizes for both user satisfaction and provider fairness. Define what ``fair exposure'' means in this context and propose a concrete optimization objective that balances the two goals.
\end{challengeexercises}
